\documentclass[openany,fontset=fandol]{ctexbook}
\usepackage{amsmath,amssymb}
\usepackage[unicode,hidelinks]{hyperref}
\setlength{\emergencystretch}{2em}
\newcommand{\Oa}{\mathcal O}
\begin{document}
\title{代数几何基础第五卷预备笔记\\奇异零点与顶次外幂判别}
\author{据英译编辑稿续译，参校法文扫描}
\date{2026年11月20日}
\maketitle
\chapter*{范围与接续}
本批从英译第71页划去段结束标记后的首句开始，连续译至第72页开头顶次外幂所给截面止。对读法文扫描第107--108页，下一起点是英译第72页接下去的微分刻画句，对应法文第109页“Pour caractériser par voie différentielle”。本批非划去正文，独立于已合并的第十七版第五卷增补。
\chapter*{零点的术语与奇异子预概形（续）}
请注意，（ii）和（ii bis）的类似条件，本可以提前放到命题9中。应当引入如下术语。设$X$为正则预概形，$E$是$X$上的局部自由模，$u:E\to\Oa_X$是同态，即$E^\vee$的一个截面。对
$x\in V(u)=V(u(E)\Oa_X)$，若$u$是$\Oa_X$-正则的，即
\[
\dim\Oa_{V(u),x}=\dim\Oa_{X,x}-\operatorname{rank}_x E,
\]
并且$x$是$V(u)$的正则点，就称$x$为$u$的\emph{非奇异零点}。取$x$的一个邻域上$E$的一组基$\phi_1,\ldots,\phi_m$，这也可以表述为：各$u(\phi_i)$是$\Oa_{X,x}$的某个正则参数系的一部分。否则，称$x$为\emph{奇异零点}。

若$x$满足推论10中的等价条件，就称$x$为\emph{普通奇异零点}（因此它确实是奇异零点）；这里自然取$A=\Oa_{X,x}$，其余记号相应解释。若$x$是奇异零点，却不是普通奇异零点，就称为\emph{超奇异零点}。因此，非超奇异零点就是非奇异零点或普通奇异零点。

最后，当$X$在域$k$上光滑时，可以用显然的方式引入这些概念的几何版本。注意，撇开维数条件，也就是$\Oa_X$-正则性条件不谈，$x$的性质实际上可从它作为$V(u)$上的点的性质看出，即可从局部环$\Oa_{V(u),x}$看出（必要时将它视为$k$-代数）。在基域为$k$的情形，若$x$是$k$-有理点，则相对于$k$的概念与绝对概念一致，这可像$E$为可逆模时那样看出。最后，若$X$在预概形$S$上光滑，而$u:E\to\Oa_X$仍为同态，$E$为有限型局部自由模，就逐纤维地定义相对于$S$的对应概念。注意，若$x$相对于$S$非超奇异，则$V(u)$在$x$处在$S$上平坦（因为$u$相对于$S$是$\Oa_X$-正则的……）。\footnote{平坦据法文第108页清楚的“plat”；英译此处作“is that”。英译者另注基底字母曾改写，可能是$Y$；本段按逐纤维讨论所用的同一基底记为$S$。}

当然（在定义正式写定以后，至少应把这点作为备注指出），也完全可以从局部自由模$E$的截面$u$出发，把$u$理解为同态$E^\vee\to\Oa_X$，即可应用前述定义。

当$X$在$Y$上光滑，给定同态$u:E\to\Oa_X$且$E$为秩$m$的局部自由模时，$u$相对于$Y$的奇异零点集是$V(u)=V^0$中一个确定的闭子预概形$V(u)^{\mathrm{sing}}=V^1$的底层集合。构造如下：取$\mathrm d^1u$，即$\mathcal P^1_{X/Y}(E^\vee)$的一个截面。它在$V^0$上的限制可以解释为同态
\[
E\otimes\Oa_{V^0}\longrightarrow
\Omega^1_{X/Y}\otimes\Oa_{V^0}.
\]
于是$V(u)^{\mathrm{sing}}=V^1$表示相应同态
\[
(\bigwedge^m E)\otimes\Oa_{V^0}\longrightarrow
\Omega^m_{X/Y}\otimes\Oa_{V^0}
\]
的零点预概形，也就是以下线丛扭曲的微分模中一个确定截面的零点预概形：
\[
\Omega^m_{X/Y}\otimes(\det E)^{-1}\otimes\Oa_{V^0}.
\]
\chapter*{译校说明（非原文）}
\paragraph{底本与基底。}
法文第107页结束标记后开头另有边栏记号，但已无此前跨页斜线。法文第108页基底处可见印字$S$与手改字样交替；英译者也标注犹疑。本稿在逐纤维定义一段统一用$S$，在下一段微分构造统一用$Y$，并不宣称每处手改字母的笔画皆可确定。地位相同的底基更名不改变公式；不能同时把不同基底的相对微分混为一个模。

\paragraph{外幂的类型。}
英译第71页把最后的源排作$E|_{V^0}$，法文第108页末明确在$E$前写$\bigwedge^m$。取前一同态的第$m$外幂，因此源为$\det E|_{V^0}$；$\Omega^m_{X/Y}$指$\bigwedge^m\Omega^1_{X/Y}$，本身不一定为线丛。局部选基后，该截面的分量为微分矩阵的$m$阶子式，故其零点预概形是所需的秩下降闭子预概形。

\paragraph{相对平坦的条件。}
非超奇异包括前文两类，均带纤维上的预期余维条件。在光滑环境的正则纤维局部环中，满足该条件的局部$m$个方程构成正则序列；光滑环境本身对基底平坦。在通常的局部平坦性判据适用时，逐个商去这些在纤维上正则的方程，得到零点预概形在该点相对平坦。此处是对原文简述的解释，不是从某个二次型在单一纤维上非退化、却无维数条件直接推出平坦。
\end{document}
